% Данный файл распространяется под лицензией CC BY 4.0. Текст лицензии размещён на https://creativecommons.org/licenses/by/4.0/
% (c) Кафедра системного программирования, 2026

% Компилировать XeTeX

\documentclass[a4paper]{article}

\usepackage[a4paper, top=8mm, bottom=8mm, left=8mm, right=8mm]{geometry}

\usepackage{polyglossia}
\setdefaultlanguage[babelshorthands=true]{russian}
\setotherlanguage{english}

\usepackage{fontspec}
\setmainfont{FreeSerif}
\newfontfamily{\russianfonttt}[Scale=0.7]{DejaVuSansMono}

\usepackage[tiny, compact]{titlesec}

\usepackage{titling}
\setlength{\droptitle}{-1cm}
\pretitle{\begin{center}\begin{bfseries}\Large}
\posttitle{\par\end{bfseries}\end{center}}
\preauthor{\begin{center}\normalsize}
\postauthor{\par\end{center}\vspace{-1.8cm}}

\usepackage{hyperref}
\usepackage{bookmark}
\usepackage{csquotes}
\usepackage{amsmath}
\usepackage{booktabs}
\usepackage{cite}


\title{Архиватор на базе алгоритма Хаффмана}

\author{Овечкин Андрей Максимович}

\date{}

\begin{document}

\maketitle

\begin{flushright}
    Группа: \emph{25.Б72-мм}

    Кафедра: \emph{Системного программирования}

    Научный руководитель: \emph{Литвинов Юрий Викторович}
    
    Номер семестра практики: \emph{2}
\end{flushright} 

\section{Введение}
Сжатие данных без потерь~--- одна из фундаментальных задач информатики. 
Алгоритм Хаффмана~\cite{huffman1952}, предложенный Дэвидом Хаффманом в 1952~году, строит 
оптимальный префиксный код для заданного распределения символов и остаётся 
основой многих современных форматов сжатия (Deflate в gzip и PNG, 
BWT~+~Huffman в bzip2). 

Реализация этого алгоритма позволяет на практике освоить работу с бинарными 
деревьями, побитовыми операциями и проектирование двоичных форматов. 
Результатом является консольный архиватор, 
написанный на языке Rust, который сжимает произвольные бинарные файлы 
и восстанавливает их без потерь.
\section{Постановка задачи}

\textbf{Целью работы является} разработка консольного архиватора на языке 
Rust, реализующего сжатие без потерь с помощью байтового кодирования 
Хаффмана с каноническими кодами и обеспечивающего точное восстановление 
исходных данных.

\textbf{Для достижения цели требуется решить следующие задачи:}
\begin{enumerate}
    \item Выполнить обзор предметной области: алгоритмы сжатия без потерь, 
          существующие архиваторы, язык Rust.
    \item Спроектировать решение: формат архива, архитектуру модулей, 
          интерфейс командной строки.
    \item Реализовать сжатие и распаковку.
    \item Провести тестирование: интеграционные тесты, краевые случаи, 
          обработка повреждённых архивов.
\end{enumerate}

\section{Обзор}

\subsection{Алгоритм Хаффмана}

Алгоритм Хаффмана решает
задачу построения оптимального префиксного кода для заданного распределения
символов. Оптимальность понимается в смысле минимальной средневзвешенной
длины кодового слова при условии, что код является префиксным (ни один код
не является префиксом другого), что позволяет однозначно декодировать
последовательность без разделителей.

Алгоритм является жадным и работает следующим образом:
\begin{enumerate}
    \item Каждый символ исходного алфавита становится листом дерева,
          весом которого является частота символа.
    \item Два узла с наименьшими весами объединяются в новый внутренний
          узел, вес которого равен сумме весов потомков.
    \item Шаг~2 повторяется, пока не останется один корневой узел.
\end{enumerate}

Ветви дерева кодируются битами: левый потомок~--- 0, правый~---~1.
Кодом символа является последовательность битов на пути от корня до листа.
Чем чаще встречается символ, тем короче его код.
Классический алгоритм является двухпроходным: первый проход по файлу
собирает статистику частот, второй~--- выполняет кодирование.

\subsection{Адаптивный алгоритм Хаффмана}

Существует также однопроходная модификация~--- адаптивный алгоритм
Хаффмана, предложенный Галлагером и Кнутом. В нём дерево строится
и обновляется на лету по мере чтения данных, что устраняет необходимость
в предварительном подсчёте частот и хранении таблицы кодов в заголовке.
Однако адаптивный вариант сложнее в реализации из-за необходимости
поддерживать инвариант дерева при каждом добавлении символа, а его степень
сжатия незначительно ниже из-за использования локальных, а не глобальных
частот.

\subsection{Канонические коды}

Прямая передача дерева Хаффмана в заголовке архива требует записи его
топологии либо полной таблицы кодов, что занимает значительный объём.
Эту проблему решают канонические коды: достаточно сохранить только длины
кодов для каждого символа, а сами коды восстанавливаются по единому
правилу. Пары \enquote{(длина, символ)} сортируются по возрастанию длины,
затем символа. Первый код принимается равным нулю, а каждый следующий
вычисляется по формуле:

\[
\text{code}_{i} = (\text{code}_{i-1} + 1) \ll (\text{length}_i - \text{length}_{i-1})
\]

\subsection{Выводы из обзора}

На основе анализа алгоритма были приняты следующие проектные решения:
\begin{itemize}
    \item выбран классический двухпроходный алгоритм Хаффмана как
          наиболее наглядный для учебной реализации;
    \item применены канонические коды для минимизации заголовка архива;
    \item реализация выполнена на языке Rust;
\end{itemize}

\section{Описание предлагаемого решения}


\subsection{Архитектура}

Проект состоит из пяти модулей:
\begin{itemize}
    \item \texttt{main.rs}~--- CLI: парсинг аргументов командной строки,
          запуск \texttt{zip} или \texttt{unzip};
    \item \texttt{lib.rs}~--- публичное API: реэкспорт функций
          \texttt{zip} и \texttt{unzip};
    \item \texttt{tree\_lib.rs}~--- структура бинарного дерева
          \texttt{Node<T,U>} и обход DFS;
    \item \texttt{huffman\_utils.rs}~--- построение дерева Хаффмана,
          генерация канонических кодов, восстановление дерева по таблице;
    \item \texttt{zip\_lib.rs}~--- логика сжатия и распаковки.
\end{itemize}

\subsection{Формат архива}
Заголовок архива имеет следующую структуру:

\begin{table}[h]
\centering
\caption{Формат заголовка архива}
\begin{tabular}{lll}
\toprule
\textbf{Смещение} & \textbf{Размер} & \textbf{Описание} \\
\midrule
0  & 8~байт & Размер исходного файла (little-endian \texttt{u64}) \\
8  & 1~байт & Количество уникальных байтов минус один \\
9  & переменный & Таблица длин кодов \\
\ldots & \ldots & Сжатые данные (побитовая упаковка) \\
\bottomrule
\end{tabular}
\end{table}

Таблица длин кодов хранится в одном из двух форматов: если количество 
уникальных символов меньше 128, каждый элемент~--- пара \texttt{(байт, длина)} 
(компактный режим). Иначе записываются 256 байт, где позиция~$i$ содержит 
длину кода для символа~$i$ (значение~0 означает отсутствие символа). 
Такой подход минимизирует заголовочные данные при любом размере алфавита.
\subsection{Алгоритм сжатия}

Сжатие состоит из четырёх этапов.

\textbf{1. Подсчёт частот.} Входной файл читается побайтово, для каждого 
байта подсчитывается частота вхождений.

\textbf{2. Построение дерева Хаффмана.} 
Из узлов-листьев с частотами собирается вектор, сортируемый по частоте. 
На каждой итерации два узла с минимальной частотой сливаются в один 
внутренний узел, который вставляется обратно с помощью бинарного поиска. 
В случае одного уникального байта добавляется фиктивный узел с частотой~0.

\textbf{3. Получение канонических кодов.} 
DFS-обход собирает глубины листьев (длины кодов). Пары 
\enquote{(длина, символ)} сортируются по длине, затем по символу. 
Первый код равен~0; каждый следующий вычисляется по формуле:
\[\text{code}_{i} = (\text{code}_{i-1} + 1) \ll (\text{length}_i - \text{length}_{i-1})\]
Канонические коды позволяют хранить в заголовке только длины, а не сами 
коды, что существенно уменьшает размер заголовка.

\textbf{4. Запись заголовка и данных.} 
В файл записываются размер оригинала, количество символов, таблица длин. 
Затем исходный файл читается повторно, для каждого байта его код 
побитово упаковывается в выходной поток; неполный последний байт 
дополняется нулями.

\subsection{Алгоритм распаковки}

Распаковка читает заголовок, восстанавливает таблицу канонических кодов 
и строит дерево: для каждого кода создаётся путь в дереве (0~--- левый 
потомок, 1~--- правый), лист помечается символом. Затем сжатые данные 
читаются побитово: спуск по дереву до листа даёт очередной байт, после 
чего указатель возвращается в корень. Процесс продолжается, пока не будет 
восстановлено ровно \texttt{original\_size} байт.

\subsection{Технические детали}

\begin{itemize}
    \item \textbf{Два режима таблицы.} При~<~128 символах хранятся только 
          существующие пары (компактно), при~>=~128~--- полный массив на~256 
          значений. Переключение по одному биту 
          исключает избыточность.
    \item \textbf{Edge-case: один уникальный байт.} Алгоритм Хаффмана 
          требует как минимум двух узлов для построения дерева; при одном 
          байте искусственно добавляется нулевой узел, что гарантирует 
          корректную работу кодирования.

\end{itemize}

\subsection{Используемые технологии}

\begin{tabular}{ll}
\textbf{Язык программирования} & Rust, edition 2024 \\
\textbf{Система сборки} & Cargo \\
\textbf{CI} & GitHub Actions (cargo test, clippy, fmt) \\
\end{tabular}

\section{Тестирование}

Для проверки корректности работы написаны 10~интеграционных тестов.

\begin{table}[h]
\centering
\caption{Интеграционные тесты}
\begin{tabular}{ll}
\toprule
\textbf{Тест} & \textbf{Проверяемый сценарий} \\
\midrule
\texttt{random\_true\_test} & Round-trip: zip $\to$ unzip, совпадение с оригиналом \\
\texttt{test\_large\_data}  & Сжатие 1~МБ однородных данных \\
\texttt{test\_empty\_file}  & Ошибка при пустом файле \\
\texttt{test\_file\_one\_byte} & Один уникальный байт \\
\texttt{test\_file\_same\_bytes} & 512 одинаковых байт \\
\texttt{test\_file\_all\_bytes\_n\_times} & Все 256~байт с разной частотой \\
\texttt{test\_file\_all\_bytes} & Все 256~байт по~1~разу \\
\texttt{test\_unzip\_short\_file\_defect} & Слишком короткий архив \\
\texttt{test\_unzip\_file\_damaged\_table} & Испорченная таблица \\
\texttt{test\_unzip\_wrong\_original\_size} & Неверный \texttt{original\_size} \\
\bottomrule
\end{tabular}
\end{table}

Тесты покрывают корректное сжатие и распаковку, граничные случаи 
(пустой файл, один символ, однородные данные, полный алфавит), 
а также обработку повреждённых входных данных.

В непрерывной интеграции (GitHub~Actions) при каждом пуше выполняются:
\begin{itemize}
    \item \texttt{cargo test ---verbose}~--- запуск всех тестов;
    \item \texttt{cargo fmt ---check}~--- проверка форматирования;
    \item \texttt{cargo clippy --- -D warnings}~--- статический анализ кода.
\end{itemize}

Для демонстрации эффективности сжатия были произведены замеры 
на различных типах данных. Замеры выполнялись на Apple~M1 (8~ГБ 
RAM, macOS~26.5.1). Время замерялось как wall-clock время 
вызова программы. Для каждого сценария выполнено по~50~запусков 
с~последующим усреднением.

\begin{table}[h]
\centering
\caption{Результаты замеров сжатия (усреднено за 50~запусков)}
\begin{tabular}{lcccc}
\toprule
\textbf{Тип данных} & \textbf{Исходный} & \textbf{Сжатый} & \textbf{Коэфф.} & \textbf{Скорость (сжатие / распаковка)} \\
\midrule
\multicolumn{5}{c}{\emph{Тестовые сценарии из интеграционных тестов}} \\
\cmidrule(lr){1-5}
\texttt{random\_true\_test} & 29~Б    & 42~Б  & 0,69 & 0,003~с / 0,003~с \\
\texttt{test\_file\_one\_byte} & 1~Б & 12~Б & 0,08 & 0,003~с / 0,003~с \\
\texttt{test\_file\_same\_bytes} & 512~Б & 75~Б & 6,83 & 0,004~с / 0,004~с \\
\texttt{test\_file\_all\_bytes} & 256~Б & 521~Б & 0,49 & 0,004~с / 0,004~с \\
\texttt{test\_all\_bytes\_n\_times} & 32~896~Б & 32~145~Б & 1,02 & 0,087~с / 0,070~с \\
\texttt{test\_large\_data} & 1~048~576~Б & 131~083~Б & 8,00 & 1,163~с / 1,719~с \\
\midrule
\multicolumn{5}{c}{\emph{Дополнительные сценарии}} \\
\cmidrule(lr){1-5}
Маленький текст & 29~Б & 58~Б & 0,50 & 0,003~с / 0,003~с \\
Средний текст & 21~500~Б & 9~481~Б & 2,27 & 0,039~с / 0,042~с \\
Большой текст & 360~000~Б & 157~533~Б & 2,29 & 0,592~с / 0,656~с \\
Однородные данные (\texttt{0xAD}, 1~МБ) & 1~048~576~Б & 131~083~Б & 8,00 & 1,164~с / 1,727~с \\
Все 256~байт & 256~Б & 521~Б & 0,49 & 0,004~с / 0,004~с \\
256~байт с~частотой $i$ & 32~896~Б & 32~145~Б & 1,02 & 0,087~с / 0,070~с \\
Сильно смещённое распределение & 999~980~Б & 174~998~Б & 5,71 & 1,189~с / 1,661~с \\
Случайные байты (100~КБ) & 102~400~Б & 102~665~Б & 1,00 & 0,273~с / 0,216~с \\
\bottomrule
\end{tabular}
\end{table}

Результаты отражают теоретические свойства алгоритма Хаффмана:
наибольшая степень сжатия достигается на данных с сильным перекосом
частот (однородные~---~8,00; смещённое распределение~---~5,71).
Тексты сжимаются примерно вдвое (2,27--2,29). При равномерном
распределении символов (случайные байты, полный алфавит) сжатие
отсутствует, а маленькие файлы (~30~Б) увеличиваются, 
что ожидаемо и не является дефектом реализации. Время сжатия
и распаковки составляет доли секунды для файлов размером до~сотен
килобайт и~1--2~секунды для файлов размером~1~МБ.

\section{Заключение}

В ходе выполнения работы получены следующие результаты:

\begin{itemize}
    \item Разработан консольный архиватор на языке Rust, реализующий 
          сжатие без потерь с помощью байтового кодирования Хаффмана.
    \item Применены канонические коды, позволяющие хранить в заголовке 
          только длины кодов и восстанавливать их на стороне декодера.
    \item Спроектирован бинарный формат архива с компактным и полным 
          представлением таблицы длин.
    \item Проведено тестирование (10~интеграционных тестов, включая 
          граничные случаи и повреждённые заголовки), настроена 
          непрерывная интеграция с линтером и запуском тестов.
    \item Выполнены замеры степени сжатия на разных типах данных.
\end{itemize}

Репозиторий: \url{https://github.com/Iseyne/Archiver-Huffman}

\begin{thebibliography}{9}
\bibitem{huffman1952} Huffman, D. (1952). A Method for the Construction of Minimum-Redundancy Codes.
\textit{Proceedings of the IRE}. 40 (9): 1098--1101.
\end{thebibliography}

\end{document}